\documentclass[11pt]{article}
\newcommand{\SheetCode}{Theory 01}
\usepackage[margin=25mm]{geometry}
\usepackage[T1]{fontenc}
\usepackage{lmodern}
\DeclareUnicodeCharacter{2605}{$\star$}
\usepackage{microtype}
\usepackage{amsmath,amssymb,mathtools,bm}
\usepackage{enumitem}
\usepackage{xcolor}
\usepackage{fancyhdr}
\usepackage{hyperref}
\usepackage[most]{tcolorbox}
\definecolor{agiBlue}{HTML}{173B57}
\definecolor{agiGold}{HTML}{B7791F}
\definecolor{agiLight}{HTML}{EFF6FA}
\definecolor{agiGreen}{HTML}{184D3B}
\hypersetup{colorlinks=true,linkcolor=agiBlue,urlcolor=agiBlue}
\setlength{\parindent}{0pt}
\setlength{\parskip}{0.55em}
\setlength{\headheight}{14pt}
\setlist{nosep,leftmargin=*}
\pagestyle{fancy}
\fancyhf{}
\fancyhead[L]{\small Part of the \href{https://github.com/mmikhasenko/GIModel.jl}{Agentic Gottried--Isgur} project}
\fancyhead[R]{\SheetCode}
\fancyfoot[C]{\thepage}
\newcommand{\dd}{\,\mathrm d}
\newcommand{\ket}[1]{\lvert #1\rangle}
\newcommand{\bra}[1]{\langle #1\rvert}
\newcommand{\braket}[2]{\langle #1\mid #2\rangle}
\newcommand{\expect}[1]{\left\langle #1\right\rangle}
\newcommand{\op}[1]{\widehat{#1}}
\newcommand{\GeV}{\,\mathrm{GeV}}
\newcommand{\R}{\mathbb{R}}
\newtcolorbox{goals}{colback=agiLight,colframe=agiBlue,title=Learning outcomes}
\newtcolorbox{reading}{colback=white,colframe=agiGold,title=Book reference,
  after skip=5mm}
\newtcolorbox{paperbridge}{colback=agiGreen!7,colframe=agiGreen,
  colbacktitle=agiGreen,coltitle=white,title=Connection to the GI paper}
\newtcolorbox{problem}[1]{colback=white,colframe=agiBlue,title={#1},
  before skip=3.5mm,after skip=1mm}
\newtcolorbox{solution}{colback=black!2,colframe=black!45,title=Solution}
\newtcolorbox{quiz}{colback=white,colframe=agiGold,title=Post-solution concept check}
\newtcolorbox{checkpoint}{colback=white,colframe=agiGold,title=Interpretation checkpoint}
\newcommand{\sheettitle}[3]{%
  \begin{center}
  {\Large\bfseries #1}\par
  \vspace{0.2em}{\color{agiBlue}#2}\par
  \vspace{0.4em}\small #3
  \end{center}\vspace{0.5em}}
\newcommand{\difficulty}[1]{\textbf{Difficulty:} #1}
\newcommand{\timebox}[1]{\textbf{Suggested time:} #1}
\newcommand{\answer}{\par\smallskip\color{black!50}\textbf{Answer:} }

\begin{document}
\sheettitle{Sheet 1: From two particles to one radial wave}{The kinematic backbone of every GI state}{Prerequisite: elementary wave mechanics. \difficulty{★/★★} \quad \timebox{2--3 hours}}

\begin{goals}
Separate center-of-mass and relative motion; derive the radial equation; move
correctly between $R(r)$ and the reduced wave $u(r)$; identify the boundary
conditions that make the radial Hamiltonian self-adjoint.
\end{goals}

\begin{reading}
Selected reading in Landau--Lifshitz, \emph{Quantum Mechanics: Non-Relativistic
Theory}: Ch. III, ``Schr\"odinger's Equation,'' and Ch. V, ``Motion in a
Centrally Symmetric Field.'' This is one useful reference selection; the sheet
is self-contained and does not follow the book chapter by chapter.
\end{reading}


We use $\hbar=c=1$. Bold symbols denote three-dimensional vectors.
In position representation the momentum of particle $i$ is the operator
$\bm p_i=-i\nabla_{\bm r_i}$, where the derivative holds the other particle's
coordinate fixed; hence $\bm p_i^2=-\nabla_{\bm r_i}^2$.
Problem 1 starts in an arbitrary inertial frame and then specializes to the
center-of-mass rest frame. For Problems 2--5, $\Psi(\bm r)$ denotes the relative
wavefunction, with $r=|\bm r|$, $\Omega=(\theta,\phi)$, and
\[
\Psi(\bm r)=R_L(r)Y_{Lm}(\Omega)=\frac{u_L(r)}rY_{Lm}(\Omega),
\qquad L=0,1,\ldots,\quad m=-L,\ldots,L.
\]
Here $R_L$ is the radial wave and $u_L=rR_L$ is the reduced radial wave;
$Y_{Lm}$ is a spherical harmonic normalized by
$\int |Y_{Lm}|^2\dd\Omega=1$.

\begin{problem}{Problem 1 ★ — Separate the two-body Hamiltonian}
\textbf{Given.} Two particles have positive masses $m_1,m_2$, positions
$\bm r_1,\bm r_2$, and a central interaction $V(|\bm r_1-\bm r_2|)$.
The nonrelativistic Hamiltonian is
\[
H=\frac{\bm p_1^2}{2m_1}+\frac{\bm p_2^2}{2m_2}
 +V(|\bm r_1-\bm r_2|).
\]
Introduce $\bm R=(m_1\bm r_1+m_2\bm r_2)/(m_1+m_2)$ and
$\bm r=\bm r_1-\bm r_2$. Define their conjugate operators by
$\bm P=-i\nabla_{\bm R}$ and $\bm p=-i\nabla_{\bm r}$, with the other new
coordinate held fixed in each derivative.

\textbf{Find.} (a) Invert the coordinate change and use the chain rule to
express $\bm p_1,\bm p_2$ in terms of $\bm P,\bm p$.
(b) Show that the Hamiltonian takes the form
$H=P^2/(2M)+p^2/(2\mu)+V(r)$, identifying $M$ and $\mu$ from its coefficients.
(c) Set the total momentum $\bm P=0$ and identify the constituent momenta.
For the supplied relativistic energies $E_i=\sqrt{\bm p_i^2+m_i^2}$,
write $T_{\mathrm{GI}}=E_1+E_2$ in this frame, including the rest energies.
No potential model or numerical masses are needed.

\textbf{Hint:} Apply the chain rule to a test wavefunction
$\Phi(\bm R,\bm r)$ before squaring the momentum operators.
\end{problem}
\begin{solution}
The inverse transformation is $\bm r_1=\bm R+(m_2/M)\bm r$ and
$\bm r_2=\bm R-(m_1/M)\bm r$. The chain rule gives
$\nabla_{\bm r_1}=(m_1/M)\nabla_{\bm R}+\nabla_{\bm r}$ and
$\nabla_{\bm r_2}=(m_2/M)\nabla_{\bm R}-\nabla_{\bm r}$, so
$\bm p_1=(m_1/M)\bm P+\bm p$ and
$\bm p_2=(m_2/M)\bm P-\bm p$. Equivalently,
$\bm P=\bm p_1+\bm p_2$ and
$\bm p=(m_2\bm p_1-m_1\bm p_2)/M$. Substitution cancels the cross term and gives
\[
M=m_1+m_2,\qquad \mu=\frac{m_1m_2}{m_1+m_2},\qquad
H=\frac{P^2}{2M}+\frac{p^2}{2\mu}+V(r).
\]
In the rest frame $\bm P=0$ and $\bm p_1=-\bm p_2\equiv\bm p$. Therefore the
sum of the two relativistic free-particle energies is
\[
T_{\mathrm{GI}}(p)=\sqrt{p^2+m_1^2}+\sqrt{p^2+m_2^2}.
\]
Unlike the quadratic nonrelativistic expression, this exact expression does not
combine the two constituent masses into a single reduced mass.
\end{solution}
\begin{quiz}
The nonrelativistic relative kinetic energy is $p^2/(2\mu)$, with
$\mu=m_1m_2/(m_1+m_2)$. The exact GI term is
$T_{\mathrm{GI}}=\sqrt{p^2+m_1^2}+\sqrt{p^2+m_2^2}$. Why is there no single
$\mu$ in the second expression? \answer the reduced mass appears only after
expanding both square roots to quadratic order in $p/m_i$; it is not an
independent parameter of the exact relativistic energies.
\end{quiz}

\begin{problem}{Problem 2 ★ — Derive the reduced radial equation}
Starting from $[-\nabla^2/(2\mu)+V(r)]\Psi=E\Psi$, derive the equation obeyed by
$u_L(r)$.

\textbf{Hint:} The input is the relative Schr\"odinger equation, a specified
central $V(r)$, the reduced mass $\mu$ from Problem 1, and the separated wave
above. Obtain a one-dimensional differential eigenvalue equation for $u_L$,
including its centrifugal term; do not solve for $E$ for an unspecified $V$.
You may use the spherical-coordinate Laplacian without deriving it:
\[
\nabla^2=\frac1{r^2}\partial_r(r^2\partial_r)+\frac1{r^2}\Delta_\Omega,
\qquad
\Delta_\Omega=\frac1{\sin\theta}\partial_\theta(\sin\theta\partial_\theta)
 +\frac1{\sin^2\theta}\partial_\phi^2,
\]
with $\Delta_\Omega Y_{Lm}=-L(L+1)Y_{Lm}$. Substitute $R_L=u_L/r$
only after applying the radial derivatives.
See also \href{https://en.wikipedia.org/wiki/Del_in_cylindrical_and_spherical_coordinates}{Wikipedia: differential operators in spherical coordinates}.
\end{problem}
\begin{solution}
The spherical Laplacian gives
$\nabla^2(RY)=r^{-2}\partial_r(r^2\partial_rR)Y-L(L+1)R Y/r^2$.
With $R=u/r$, the first derivative cancels:
\[
\left[-\frac1{2\mu}\frac{d^2}{dr^2}
+\frac{L(L+1)}{2\mu r^2}+V(r)\right]u_L(r)=E u_L(r).
\]
The centrifugal term is kinematic; it is not part of the GI potential.
\end{solution}
\begin{quiz}
Suppose a numerical method represents the radial state by its values at selected
radii $r_1,r_2,\ldots$. Which function does it represent in this course?
\answer the reduced radial function $u(r)$, not $R(r)$. This matters because
$u$ is normalized with the one-dimensional integral $\int|u|^2\dd r$; the
factor $r^2$ from the three-dimensional volume element has already cancelled
against the $1/r^2$ in $|R|^2=|u|^2/r^2$.
\end{quiz}

\begin{problem}{Problem 3 ★ — Prove the normalization conversion}
Prove that $\int |\Psi|^2\dd^3r=1$ is equivalent to
$\int_0^\infty |u_L(r)|^2\dd r=1$ when $Y_{Lm}$ is normalized.

\textbf{Hint:} Use $\dd^3r=r^2\dd r\dd\Omega$ with
$\dd\Omega=\sin\theta\dd\theta\dd\phi$, $0\le\theta\le\pi$ and
$0\le\phi<2\pi$. Show the cancellation of the radial Jacobian explicitly.
\end{problem}
\begin{solution}
Use $\dd^3r=r^2\dd r\dd\Omega$ and $R=u/r$:
\[
\int|\Psi|^2\dd^3r=\int_0^\infty r^2\frac{|u|^2}{r^2}\dd r
\int|Y_{Lm}|^2\dd\Omega=\int_0^\infty|u|^2\dd r.
\]
Consequently a central expectation is $\expect{f(r)}=\int|u|^2f(r)\dd r$.
\end{solution}
\begin{quiz}
For a spherically symmetric density, the volume element gives
$\int\rho(\bm r)\dd^3r=4\pi\int r^2\rho(r)\dd r$. Should that factor
$4\pi r^2$ be inserted again into an $S$-wave expectation written with the
normalized reduced wave $u$? \answer no. The angular integral and the
$r^2|R|^2=|u|^2$ cancellation were already used to obtain
$\int|u|^2f(r)\dd r$; inserting them again would count the geometry twice.
\end{quiz}

\begin{problem}{Problem 4 ★★ — Determine the origin boundary condition}
Use the radial equation from Problem 2, with $r^2V(r)\to0$ as $r\to0$.
Find the two leading powers $u_L\sim r^a$ near the origin by substituting this
ansatz into the leading equation. Select the usual regular three-dimensional
solution, with no point interaction at the origin, and state its boundary
condition on $u_L$. Check $L=0$ separately: is square integrability alone enough
to choose the physical branch?
\end{problem}
\begin{solution}
The dominant equation is $-u''+L(L+1)u/r^2=0$. Hence
$a(a-1)=L(L+1)$, so $a=L+1$ or $a=-L$. For $L\ge1$, the second branch is not square integrable in $\dd r$.
For $L=0$, however, $u\sim\mathrm{constant}$ is locally square integrable;
it gives $R\sim1/r$ and is excluded by the regular three-dimensional domain,
not by normalization alone. With no point interaction at the origin, the
physical branch is $u_L(r)\sim r^{L+1}$ and $u_L(0)=0$.
\end{solution}
\begin{quiz}
To solve the differential equation numerically, choose radii
$r_i=ih$ separated by a step $h$. Why may the list of unknown values begin at
$r_1=h$ rather than $r_0=0$? \answer regularity has already specified the
boundary value $u(0)=0$. Such a specified value is called a Dirichlet boundary
condition. The quantities still to be determined are $u(h),u(2h),\ldots$.
\end{quiz}

\begin{problem}{Problem 5 ★★ — Connect the boundary condition to Hermiticity}
An operator $A$ is symmetric on a stated set of functions if
$\braket{u}{Av}=\braket{Au}{v}$ for every pair $u,v$ in that set. Show by
integration by parts that $A=-d^2/dr^2$ has this property for radial functions
obeying $u(0)=u(\infty)=0$. State what fails if the origin condition is omitted.

\textbf{Hint:} The inner product is
$\braket uv=\int_0^\infty u^*(r)v(r)\dd r$. Take $u,v$ and their first
derivatives sufficiently regular for two integrations by parts, with the
functions and derivatives decaying at infinity. Keep the endpoint term
explicit. Prove symmetry on this domain; a proof of self-adjointness is not
requested. Ask whether the term vanishes for arbitrary endpoint values if no
origin condition is imposed (other consistent boundary conditions are possible).
\end{problem}
\begin{solution}
For $u,v$ in the domain,
\[
\braket{u}{-v''}-\braket{-u''}{v}
=\left[u'^*v-u^*v'\right]_0^\infty=0.
\]
Without the boundary conditions, the surface term need not vanish, so the two
inner products need not be equal. With an appropriate complete specification of
its domain, this symmetric radial operator is self-adjoint. Self-adjointness is
the mathematical condition that ensures real energy eigenvalues and the usual
orthogonality and spectral interpretation of energy eigenstates.
\end{solution}
\begin{quiz}
Why is the formula $-d^2/dr^2$ alone not a complete radial kinetic-energy
operator? \answer one must also state which functions it acts on and their
boundary conditions. Different boundary conditions change the surface term and
can define different operators and spectra; multiplying an eigenfunction by a
normalization constant does not supply this missing information.
\end{quiz}

\begin{checkpoint}
You are ready for Sheet 2 if you can explain, without formulas, why a meson
spectrum calculation may store a one-dimensional wave while still describing a
three-dimensional state.
\end{checkpoint}

\begin{paperbridge}
GI Eqs. (1a)--(1b) define the rest-frame eigenvalue problem and its exact
kinetic energy
$H_0=\sqrt{p^2+m_1^2}+\sqrt{p^2+m_2^2}$. Eq. (2a) displays the
nonrelativistic expansion used in Problem 1. The angular reduction and radial
normalization are standard quantum mechanics needed to represent the
$\ket{jm;ls}$ sectors named below GI Eq. (14).
\end{paperbridge}
\end{document}
